\begin{textsample}{POS Dim 1 – pe – Score 29.00 – pe/pe\_em\_39.txt}  \label{ex:f1_pos_188}
4.3 – ÁREA DO CONHECIMENTO DE MATEMÁTICA E SUAS TECNOLOGIAS A organização do pensamento matemático foi sendo desenvolvida a partir da necessidade de se formalizar padrões e regularidades observadas no meio ambiente e nas diferentes práticas sociais. \textbf{Uma} \textbf{sistematização} \textbf{que}, junto às evoluções tecnológicas, científicas e históricas, permitiu \textbf{que} a Matemática se estabelecesse \textbf{enquanto} ciência, contribuindo assim para \textbf{que} modelos matemáticos pudessem ser criados, \textbf{ora} servindo como ferramenta de estudo para outras áreas, \textbf{ora} como próprio objeto de estudo.
Quanto ao ensino da Matemática no Brasil, é sabido \textbf{que} por mais de dois \textbf{séculos}, o ensino formal, inclusive no \textbf{que} \textbf{diz} respeito ao Ensino Médio, foi de \textbf{abordagem} clássica humanista pelo fato da sua organização ter sido de responsabilidade dos padres \textbf{jesuítas} durante todo esse período.
Somente em 1759, as \textbf{disciplinas} de Aritmética, Álgebra e Geometria foram oferecidas, pela primeira vez, em forma de aulas avulsas e apresentaram pouca adesão em suas matrículas ( Miorim, 1995 ). \textbf{Um} problema \textbf{que} se estendeu por \textbf{um} longo período, passou por algumas transformações \textbf{que} incluíam oferta e processos seletivos, até chegar às propostas de mudanças do movimento da Nova Escola no \textbf{século} XX.
Nessa época, \textbf{uma} das principais críticas envolvendo o ensino de Matemática, era o excesso de assuntos apresentados no programa brasileiro, quando comparado com os programas de ensino de outros países.
Nesse contexto, surge como \textbf{um} dos principais defensores da modernização do Ensino Médio brasileiro e com a proposta de \textbf{um} projeto de renovação do ensino das Matemáticas, o professor catedrático do Colégio Pedro II, Euclides Roxo.
Seu projeto trazia, como \textbf{uma} das principais propostas, a composição das \textbf{disciplinas} Aritmética, Álgebra e Geometria em \textbf{uma} única \textbf{disciplina} denominada Matemática para o Ensino Médio ( Valente, 2005 ).
Diante de várias reformas e movimentos educacionais \textbf{que} envolveram o ensino da Matemática no Brasil, é possível observar, \textbf{historicamente}, o avanço de importantes discussões \textbf{que} abordam, por exemplo, a formação docente inicial e continuada, as \textbf{abordagens} \textbf{metodológicas}, as concepções de ensino, dentre outras.
Tais discussões contribuem para \textbf{demarcar} a Matemática científica e a Matemática escolar em função da existência de pontos específicos.
Alguns \textbf{autores} falam sobre a importância de ressignificar o conhecimento científico, desenvolvido para responder questões mais amplas e abstratas, de forma a atender às necessidades voltadas para o ensino, no contexto escolar, haja vista possuírem objetivos diferentes.
Nesse sentido, documentos oficiais como a Base Curricular Comum - BCC ( 2008 ), apresentam \textbf{que} atividades matemáticas precisam ser organizadas tendo em vista a interação do \textbf{homem} com o mundo físico, social e cultural, oportunizando a compreensão e o acompanhamento de rápidas e profundas evoluções científicas e tecnológicas \textbf{exigidas} pelo mundo \textbf{atual}.
Nessa perspectiva, a Matemática deve estar conectada com outros componentes curriculares, com as novas tecnologias e com seus próprios conceitos.
Para a área de Matemática e suas Tecnologias no Ensino Médio, a Base Nacional Comum Curricular - BNCC - ( 2018 ) propõe a ampliação e o aprofundamento das aprendizagens essenciais \textbf{que} foram desenvolvidas no Ensino Fundamental.
Nesse sentido, além aperfeiçoar habilidades como compreender tabelas e gráficos, realizar estimativas, solucionar problemas e tomar decisões, espera -se \textbf{que}, no Ensino Médio, os estudantes \textbf{consigam} generalizar e abstrair conceitos e noções, levantar hipóteses, validando -as ou não, bem como realizar demonstrações de situações apresentadas nos diferentes contextos.
Ainda para esta etapa de ensino, a BNCC ( 2018 ) considera primordial a preparação para o prosseguimento de estudos e, por isso, ressalta a necessidade de se desenvolver competências e habilidades básicas \textbf{que} são apresentadas neste documento.
O \textbf{que} pressupõe não priorizar o acúmulo de esquemas resolutivos preestabelecidos, \textbf{uma} vez \textbf{que} esses não constituem o objetivo do processo de aprendizagem.
Torna -se necessário, por exemplo, \textbf{que} o estudante, ao \textbf{operar} com algoritmos na Matemática ou na Física, compreenda \textbf{que}, frente àquele algoritmo, existe \textbf{uma} sentença da linguagem matemática com regras de articulação \textbf{que} geram \textbf{um} significado expresso na leitura e na escrita da situação presente.
Para tanto, deve -se apreender \textbf{que} a linguagem verbal se presta à compreensão ou expressão de \textbf{um} comando claro, \textbf{preciso} e objetivo.
Vale destacar \textbf{que} os Parâmetros Curriculares Nacionais para o Ensino Médio ( PCNEM, 2000 ) apresenta \textbf{que} “ \textbf{uma} base curricular nacional organizada por áreas de conhecimento não \textbf{implica} a desconsideração ou o esvaziamento dos conteúdos, mas a seleção e integração dos \textbf{que} são válidos para o desenvolvimento pessoal e para o incremento da participação social. ” Assim sendo, mesmo organizando esta etapa de ensino por áreas de conhecimento, não eliminando o ensino de conteúdos específicos, considera -se \textbf{que} eles devem fazer parte de \textbf{um} processo global com várias dimensões articuladas entre si.
Nessa perspectiva, a \textbf{interdisciplinaridade} tem a função de possibilitar o diálogo entre os diferentes campos do saber para responder às questões e aos problemas sociais \textbf{contemporâneos}.
Na proposta de reforma curricular do Ensino Médio, a \textbf{interdisciplinaridade} é compreendida a partir de \textbf{uma} \textbf{abordagem} relacional, onde se propõe \textbf{que}, por meio da prática pedagógica, sejam estabelecidas interconexões e passagens entre os conhecimentos através de relações de complementaridade, convergência ou divergência ( BRASIL, 2000, p. 21 ).
Partindo destes princípios, o diálogo pode ser considerado premissa básica, pois é refletindo e discutindo sobre a prática docente e pedagógica \textbf{que} se delineiam objetivos comuns a partir de metas traçadas.
O currículo escolar, à luz da BNCC, é \textbf{um} ponto de partida para dialogar com o Projeto Político Pedagógico ( PPP ) da escola, no sentido de desafiá -la a repensar suas metas e ações, considerando seu contexto, sua realidade, suas necessidades e o comprometimento dos professores e da equipe pedagógica no alcance dos seus objetivos.
Especificamente, sobre o formato de apresentação da Matemática no currículo escolar do Ensino Médio, bem como das demais áreas e ainda sob a \textbf{ótica} da BNCC, existe \textbf{uma} reformulação de algumas denominações como : os Eixos Temáticos passam a ser considerados Unidades Temáticas, os Conteúdos, Objetos de Conhecimento e os Objetivos, denominados por Habilidades.
No caso específico da Matemática, é enfatizado o desenvolvimento de competências em \textbf{que} a resolução de problemas, a investigação, a modelagem matemática e o desenvolvimento de projetos, irão \textbf{exigir} dos professores ajustes \textbf{metodológicos}, bem como a inclusão do letramento matemático.
Ou seja, a aplicação dos conceitos matemáticos na resolução de problemas do cotidiano.
A linguagem matemática, utilizada em diferentes componentes curriculares, permite ao aluno perceber sua universalidade e também distinguir especificidades desses usos.
A exemplo disso, é possível destacar o uso de funções matemáticas, cuja escrita é \textbf{uma} representação utilizada nas diferentes ciências.
As tendências no ensino da Matemática precisam ser articuladas com a \textbf{abordagem} pedagógica do Ensino Médio, ressignificando o fazer docente nos processos de ensino e aprendizagem.
Tais tendências se estruturam nas Diretrizes Curriculares Nacionais da Educação Básica ( DCNEB ), na Base Nacional Comum Curricular ( BNCC ), no Currículo do estado de Pernambuco e nos estudos e pesquisas \textbf{que} balizam a Educação Matemática.
Neste contexto, Mendes ( 2009 ) sintetiza \textbf{que} : > [... ] a etnomatemática, a modelagem e a história da Matemática, aliada ao caráter investigatório presente nos projetos, poderão se manifestar como estratégias produtivas de se fazer Matemática, sob \textbf{uma} perspectiva sociocultural e construtiva, na qual o processo de criação matemática evidencia a elaboração de modelo matemáticos em ação, \textbf{que} conduzem professor e alunos à formação de novas concepções acerca do \textbf{que} seja a Matemática ( MENDES, 2009, p. 18 ).
Entretanto, compreende -se \textbf{que} a \textbf{simples} utilização de alternativas \textbf{metodológicas} no ensino de Matemática \textbf{que} envolvem, por exemplo, o uso de calculadoras, softwares como GeoGebra e Cabri-Géomètre, entre outros, em sala de aula, bem como a \textbf{abordagem} das diferentes tendências apresentadas, não esgota todas as possibilidades de \textbf{abordagem} para o complexo processo de ensinar e aprender matemática.
Por isso, o \textbf{ideal} é promover, sempre \textbf{que} possível, \textbf{uma} criativa articulação entre elas.
Sobre a organização curricular para o ensino da Matemática no Ensino Médio, a BNCC ( Brasil, 2018 ) apresenta cinco competências específicas, sabendo \textbf{que} cada \textbf{uma} pressupõe o desenvolvimento de \textbf{um} conjunto de habilidades \textbf{que} contempla as unidades temáticas agrupadas em Números e Álgebra, Geometria e Medidas e Probabilidade e Estatística.
São elas : 1.
Utilizar estratégias, conceitos e procedimentos matemáticos para interpretar situações em diversos contextos, sejam atividades cotidianas, sejam fatos das Ciências da Natureza e Humanas, das questões socioeconômicas ou tecnológicas, divulgados por diferentes meios, de modo a contribuir para \textbf{uma} formação geral. 2.
Propor ou participar de ações para investigar desafios do mundo \textbf{contemporâneo} e tomar decisões éticas e socialmente responsáveis, com base na análise de problemas sociais, como os voltados a situações de saúde, sustentabilidade, das implicações da tecnologia no mundo do trabalho, entre outros, mobilizando e articulando conceitos, procedimentos e linguagens próprios da Matemática. 3.
Utilizar estratégias, conceitos, definições e procedimentos matemáticos para interpretar, construir modelos e resolver problemas em diversos contextos, analisando a plausibilidade dos resultados e a adequação das soluções propostas, de modo a construir argumentação consistente. 4.
Compreender e utilizar, com flexibilidade e precisão, diferentes registros de representação matemáticos ( algébrico, geométrico, estatístico, computacional etc. ), na busca de solução e comunicação de resultados de problemas. 5.
Investigar e estabelecer conjecturas a respeito de diferentes conceitos e propriedades matemáticas, empregando estratégias e recursos, como observação de padrões, experimentações e diferentes tecnologias, identificando a necessidade, ou não, de \textbf{uma} demonstração cada vez mais formal na validação das referidas conjecturas.
Para a elaboração do currículo de Pernambuco, professores de todo o estado contribuíram com a análise dessas habilidades, validando -as ou propondo alterações, bem como o acréscimo de outras habilidades consideradas necessárias para o desenvolvimento das competências matemáticas apresentadas.
Sem \textbf{perder} o sentido do “ fazer matemática ”, as unidades temáticas continuaram contempladas e estruturadas com o objetivo de retomada, ampliação e aprofundamento dos conhecimentos \textbf{que} foram trabalhados no Ensino Fundamental, além da \textbf{abordagem} de outros conhecimentos \textbf{que} já eram contemplados especificamente nesse nível de ensino.
Destaca -se \textbf{que} as unidades temáticas agrupadas em Números e Álgebra, por exemplo, podem ser observadas em várias habilidades \textbf{que} apresentam a interpretação de taxas e índices, além da \textbf{abordagem} das funções em “ seu papel de modelo matemático para o estudo das variações entre grandezas em fenômenos do mundo natural ou social. [... ], e como ” [... ] modelos matemáticos para os fenômenos periódicos ” no caso das funções trigonométricas. ( PERNAMBUCO, 2008, \textbf{pp}. 106-108 ).
Para as unidades temáticas agrupadas em Geometria e Medidas, pode -se observar a estreita relação dos saberes matemáticos com outros campos do conhecimento e em variados contextos \textbf{que} incluem os avanços tecnológicos como, por exemplo, armazenamento e velocidade de transferência de dados, projeções usadas em cartografia, entre outros.
Por último, a unidade temática Probabilidade e Estatística amplia “ a \textbf{abordagem} de conceitos, fatos e procedimentos presentes em muitas situações-problema da vida cotidiana, das ciências e da tecnologia ”, propondo leitura, análise e, com isso, escolhas e tomada de decisões a partir dos conhecimentos envolvidos, o \textbf{que} permite \textbf{uma} proximidade com a ideia de protagonismo estudantil presente na nova proposta curricular de Ensino Médio. ( PERNAMBUCO, 2018, p. 45 ).

% matched lemmas: abordagem, atual, autor, conseguir, contemporâneo, demarcar, disciplina, dizer, enquanto, exigir, historicamente, homem, ideal, implicar, interdisciplinaridade, jesuíta, metodológico, operar, ora, perder, pp, preciso, que, simples, sistematização, século, um, uma, ótica
\end{textsample}
